% =====================================================================
%  TEMPLATE: nli-dense  —  a TIGHTER one-page variant
%
%  WHEN TO USE THIS INSTEAD OF plain-professional:
%  Two cases. (1) Your one-pager reads sparse or top-heavy: this
%  variant buys vertical room. (2) Your history genuinely needs two
%  pages: this format PERMITS 1 OR 2, where plain-professional is a
%  strict one-pager.
%
%  It is NOT "better". It trades breathing room for density.
%
%  Derived from the structure of a widely circulated PM resume template.
%
%  ⚠️ PACKAGE GUARDS, and they matter. This variant wants `ragged2e`
%  and `needspace`, which a minimal TinyTeX may not carry. Both are
%  loaded ONLY if present, each with a plain-LaTeX fallback, so the
%  file compiles either way. Do not "simplify" the guards away: a
%  template that fails to compile on the recipient's install is worse
%  than one that renders slightly plainer.
% =====================================================================
% PAGE-LIMIT: 1-2   % this format permits one or two pages; verify_resume reads this line
\documentclass[10pt,letterpaper]{article}

\usepackage[left=0.75in,right=0.75in,top=0.55in,bottom=0.55in]{geometry}
\usepackage[T1]{fontenc}
\usepackage[scaled]{helvet}
\renewcommand{\familydefault}{\sfdefault}
\usepackage{titlesec}
\usepackage{hyperref}
\hypersetup{colorlinks=true, urlcolor=black, linkcolor=black, pdftitle={[YOUR_NAME] - CV}}

% Ragged right (no hyphenation) if ragged2e is installed; plain \raggedright otherwise.
\IfFileExists{ragged2e.sty}{%
  \usepackage[document]{ragged2e}%
  \setlength{\RaggedRightRightskip}{0pt plus 2.5cm}%
}{%
  \raggedright%
}
% Keeps a role block from splitting across the page. No-op if needspace is absent.
\IfFileExists{needspace.sty}{\usepackage{needspace}}{\providecommand{\needspace}[1]{}}

\setlength{\parindent}{0pt}
% ⭐ THE `plus` IS THE POINT. Stretchable glue lets LaTeX distribute leftover
% vertical space across every paragraph break, so a SHORT history opens up on
% its own instead of stacking at the top and leaving a gap at the bottom.
% Do not replace it with a flat length, and do not hand-tune margins per resume:
% that makes every build a one-off and you lose "they all look the same".
\setlength{\parskip}{3pt plus 1pt}

\titleformat{\section}{\bfseries\fontsize{11}{13}\selectfont}{}{0em}{}[\vspace{-4pt}\hrule height 0.6pt]
\titlespacing*{\section}{0pt}{6pt}{3pt}

\let\olditemize\itemize \let\endolditemize\enditemize
\renewenvironment{itemize}
  {\begin{olditemize}\setlength{\itemsep}{1pt}\setlength{\parskip}{0pt}\setlength{\topsep}{0pt}\setlength{\leftmargini}{1.1em}}
  {\end{olditemize}}

% A role heading that will not be left stranded at the foot of the page.
\newcommand{\subblock}[1]{\needspace{5\baselineskip}\vspace{1pt}\textit{\textbf{#1}}\vspace{-6pt}}

% ⛔ THE UNFILLED SKELETON MUST COMPILE. The [PLACEHOLDER] tokens carry underscores, and a bare `_`
% is a math subscript in LaTeX, so this file failed with ~95 errors before a single field was
% filled in — while install.sh tells you to compile it as your smoke test and /add-template
% requires a test compile before registering. `underscore` makes `_` behave as text; the catcode
% fallback does the same on an install that lacks the package. Harmless once your real content is
% in: a resume has no math.
\IfFileExists{underscore.sty}{\usepackage{underscore}}{\catcode`\_=12\relax}

\begin{document}

% ============================================================
%     HEADER  — the title line under the name is this variant's tell
% ============================================================
{\fontsize{20}{22}\selectfont\bfseries [YOUR_NAME]}\\[2pt]
{\fontsize{11}{13}\selectfont [YOUR_TARGET_TITLE]}\\[3pt]
[YOUR_LOCATION] $\cdot$ [YOUR_REMOTE_STATUS] $\mid$ [YOUR_PHONE] $\mid$ [YOUR_EMAIL] $\mid$ [YOUR_LINKEDIN] $\mid$ [YOUR_WEBSITE]

% ============================================================
%     SUMMARY  (hard cap: 300 characters)
% ============================================================
\section*{Summary}
[PROFILE_SUMMARY_MAX_300_CHARS]

% ============================================================
%     CORE SKILLS
% ============================================================
\section*{Core Skills}
\textbf{[SKILL_CATEGORY_1]:} [skills, comma separated]\\
\textbf{[SKILL_CATEGORY_2]:} [skills, comma separated]\\
\textbf{[SKILL_CATEGORY_3]:} [skills, comma separated]

% ============================================================
%     EXPERIENCE  — reverse chronological, ALWAYS. Tailor by rewriting
%     the summary and bullets, NEVER by reordering the timeline.
% ============================================================
\section*{Experience}

% ⭐ THE ACHIEVEMENTS / RESPONSIBILITIES SPLIT IS THE SUBSTANTIVE HALF OF THIS FORMAT, and until
% 2026-08-10 this skeleton shipped without it (kit issue #19). `\subblock` was defined above and
% never used, so the template carried the LAYOUT (margins, parskip, needspace) and none of the
% STRUCTURE, while its name promised the format. An operator following it believed they were using
% the format and were using the margins.
%
% 📊 WHY IT IS WORTH THE LINES, measured by the partner who reported it: restructuring identical
% content into this shape took a page from 69.8% full to about 83%, because sub-headings consume
% space structurally instead of by padding. Nothing was invented and no line was added.
%
% ⚖️ ACHIEVEMENTS FIRST, and keep it shorter than Responsibilities. Achievements are outcomes with
% a number attached; Responsibilities are scope, and follow the formula ACTION + KEYWORDS + SCOPE +
% IMPACT. If a role has no achievement worth a number, that is worth knowing before the résumé goes
% out — do not pad the section to fill it.
\noindent\textbf{[COMPANY_1]} \hfill [LOCATION_1] $\cdot$ [DATES_1]\\
\textit{[TITLE_1]} \hfill \textit{[DOMAIN_TAG_1]}
% \item{} — the empty group stops LaTeX reading a [BRACKETED] placeholder as the
    % item's optional LABEL, which silently renders the bullet with no marker.
\subblock{Achievements}
\begin{itemize}
    \item{} [ACHIEVEMENT_1]
    \item{} [ACHIEVEMENT_2]
\end{itemize}
\subblock{Responsibilities}
\begin{itemize}
    \item{} [BULLET_1]
    \item{} [BULLET_2]
    \item{} [BULLET_3]
\end{itemize}

\noindent\textbf{[COMPANY_2]} \hfill [LOCATION_2] $\cdot$ [DATES_2]\\
\textit{[TITLE_2]} \hfill \textit{[DOMAIN_TAG_2]}
\subblock{Achievements}
\begin{itemize}
    \item{} [ACHIEVEMENT_1]
\end{itemize}
\subblock{Responsibilities}
\begin{itemize}
    \item{} [BULLET_1]
    \item{} [BULLET_2]
\end{itemize}

% ============================================================
%     EDUCATION & CERTIFICATIONS
% ============================================================
\section*{Education \& Certifications}
% \newline, not \\: a line break followed by a [BRACKETED] placeholder makes LaTeX read the
% placeholder as \\'s optional vertical-space argument, and the skeleton fails to compile.
\textbf{[SCHOOL]} --- [DEGREE]\newline
[CERTIFICATION], [ISSUING_BODY], [YEAR]

\end{document}
